% BEGIN R28_REGIONAL_CLOSURES
% Inserción después de REG01. No es un documento autónomo.
% Procedencia: U012_biografia_estructural_pi.tex, líneas 175--324.
% Testigos orientados: verificar_cierres_regionales.py; mismo problema,
% representantes distintos de los paseos impresos en U012.
% Emisor SHA-256: 42fdacab582de8e2d2cb7dfafc6204d200aef01249b3dcc37949e39507f066d8.
\subsubsection{Multigrafo de bloques y cinco cierres regionales mínimos}
\label{regcl27:seccion}

Las \(468\) emisiones del conjunto \(\mathcal U_6\), producidas por
el emisor APP--TRIT--TPK, determinan un multigrafo de bloques. La
identificación de sus anclajes regionales permite resolver un problema
de cobertura conjunta: recorrer las incidencias APP, de propagación,
de autoescala y de cada región de clausura en un mismo paseo cerrado.
El ciclo local de seis empalmes de la sección~\ref{sec:rec27-empalmes}
y estas coberturas conservan sus respectivas condiciones de recorrido.

\paragraph{Grafo de bloques y cociente exacto de fase.}
Para \(U=(U_1,\ldots,U_6)\in\mathcal U_6\), definimos
\begin{equation}
 A(U)=(U_1,U_2,U_3),\qquad B(U)=(U_4,U_5,U_6).
 \label{regcl27:mitades}
\end{equation}
El multigrafo no dirigido \(\widetilde G\) tiene por vértices las ternas
que aparecen como \(A(U)\) o \(B(U)\), y una arista etiquetada por cada
\(U\), incidente en esas dos ternas. Se retienen las \(468\) etiquetas:
dos emisiones distintas no se identifican por compartir extremos, ni
porque una intercambie las mitades de otra.

La fase interna actúa sobre cada terna por rotación cíclica:
\begin{equation}
 (a,b,c)\sim(b,c,a)\sim(c,a,b),\qquad
 \operatorname{can}(a,b,c)
 =\min_{\rm lex}\{(a,b,c),(b,c,a),(c,a,b)\}.
 \label{regcl27:fase}
\end{equation}
El grafo \(G\) sustituye los extremos de la arista \(U\) por
\(\operatorname{can}A(U)\) y \(\operatorname{can}B(U)\), conservando
su etiqueta. La equivalencia actúa exactamente por las tres rotaciones
cíclicas de cada vértice; cada emisión mantiene su identidad como arista.

La enumeración del mismo emisor y la búsqueda de componentes dan
\begin{equation}
\begin{array}{c|r|l}
 & |V|&\text{tamaños de componentes}\\ \hline
 \widetilde G&96&(28,28,28,4,4,4)\\
 G&32&(28,4)
\end{array}
\label{regcl27:componentes}
\end{equation}
Para obtener estos datos se reúnen los extremos distintos, se añade
cada incidencia en ambos sentidos y se recorre por anchura cada
componente aún no visitada. La adyacencia puede deduplicarse para contar
componentes, pero las etiquetas paralelas se conservan al buscar cierres
con aristas prescritas.

\paragraph{Anclaje APP y numeración regional.}
Las tres aristas comunes son
\begin{equation}
\begin{aligned}
 U_{\rm APP}&=903|850|850|252|109|252,\\
 U_e&=850|963|890|149|252|212,\\
 U_\varphi&=830|943|830|109|252|252.
\end{aligned}
\label{regcl27:anclajes}
\end{equation}
Fijamos el extremo
\(v_{\rm APP}=\operatorname{can}A(U_{\rm APP})=850|850|903\).
En las cinco regiones ya construidas se adopta el siguiente orden,
que sitúa la región transversal en la quinta fila:
\begin{equation}
\begin{array}{c|l|c|r}
 i&U_\pi^{(i)}&\rho_i&m_i\\ \hline
 1&810|923|870|169|272|272&++&144\\
 2&810|983|810|169|272|272&++&144\\
 3&870|923|810|169|272|272&++&144\\
 4&870|923|810|418|674|521&--&144\\
 5&501|614|498|169|272|272&\perp&432
\end{array}
\label{regcl27:regiones}
\end{equation}
La numeración distingue las tres regiones \(++\), el retorno mutado
\( -- \) y la transversal \(\perp\). Por ello,
\(\rho=(++ ,++ ,++ ,--,\perp)\) y
\(m=(144,144,144,144,432)\) en este orden; la suma sigue siendo \(1008\).

\paragraph{Cinco testigos completos con sentido de recorrido.}
En cada tabla, la fila \(j\) significa recorrer la arista etiquetada
\(e_j\) desde \(v_{j-1}\) hasta \(v_j\). Los extremos son representantes
canónicos de fase. Como \(G\) es no dirigido, la pareja
\((v_{j-1},v_j)\) fija el sentido de recorrido de la arista etiquetada.
Cada tabla contiene exactamente ocho aristas y satisface
\(v_0=v_8=v_{\rm APP}\).

La búsqueda por anchura descrita en la prueba proporciona los siguientes
testigos, uno para cada conjunto de cuatro aristas objetivo. La
minimalidad se refiere a la longitud; admite varios paseos que la realizan.
\paragraph{Testigo \(i=1\).}
\begin{center}
\small
\begin{tabular}{rlll}
\hline
\(j\)&\(v_{j-1}\)&\(e_j\)&\(v_j\)\\
\hline
1&\(850|850|903\)&\(109|252|252|850|903|850\)&\(109|252|252\)\\
2&\(109|252|252\)&\(109|252|252|810|923|870\)&\(810|923|870\)\\
3&\(810|923|870\)&\(810|923|870|169|272|272\)&\(169|272|272\)\\
4&\(169|272|272\)&\(169|272|272|850|963|890\)&\(850|963|890\)\\
5&\(850|963|890\)&\(850|963|890|149|252|212\)&\(149|252|212\)\\
6&\(149|252|212\)&\(149|252|212|830|943|830\)&\(830|830|943\)\\
7&\(830|830|943\)&\(830|943|830|109|252|252\)&\(109|252|252\)\\
8&\(109|252|252\)&\(903|850|850|252|109|252\)&\(850|850|903\)\\
\hline
\end{tabular}
\end{center}

\paragraph{Testigo \(i=2\).}
\begin{center}
\small
\begin{tabular}{rlll}
\hline
\(j\)&\(v_{j-1}\)&\(e_j\)&\(v_j\)\\
\hline
1&\(850|850|903\)&\(169|272|272|850|903|850\)&\(169|272|272\)\\
2&\(169|272|272\)&\(169|272|272|810|983|810\)&\(810|810|983\)\\
3&\(810|810|983\)&\(810|983|810|169|272|272\)&\(169|272|272\)\\
4&\(169|272|272\)&\(169|272|272|850|963|890\)&\(850|963|890\)\\
5&\(850|963|890\)&\(850|963|890|149|252|212\)&\(149|252|212\)\\
6&\(149|252|212\)&\(149|252|212|830|943|830\)&\(830|830|943\)\\
7&\(830|830|943\)&\(830|943|830|109|252|252\)&\(109|252|252\)\\
8&\(109|252|252\)&\(903|850|850|252|109|252\)&\(850|850|903\)\\
\hline
\end{tabular}
\end{center}

\paragraph{Testigo \(i=3\).}
\begin{center}
\small
\begin{tabular}{rlll}
\hline
\(j\)&\(v_{j-1}\)&\(e_j\)&\(v_j\)\\
\hline
1&\(850|850|903\)&\(169|272|272|850|903|850\)&\(169|272|272\)\\
2&\(169|272|272\)&\(169|272|272|850|963|890\)&\(850|963|890\)\\
3&\(850|963|890\)&\(850|963|890|149|252|212\)&\(149|252|212\)\\
4&\(149|252|212\)&\(149|252|212|870|923|810\)&\(810|870|923\)\\
5&\(810|870|923\)&\(870|923|810|169|272|272\)&\(169|272|272\)\\
6&\(169|272|272\)&\(169|272|272|830|943|830\)&\(830|830|943\)\\
7&\(830|830|943\)&\(830|943|830|109|252|252\)&\(109|252|252\)\\
8&\(109|252|252\)&\(903|850|850|252|109|252\)&\(850|850|903\)\\
\hline
\end{tabular}
\end{center}

\paragraph{Testigo \(i=4\).}
\begin{center}
\small
\begin{tabular}{rlll}
\hline
\(j\)&\(v_{j-1}\)&\(e_j\)&\(v_j\)\\
\hline
1&\(850|850|903\)&\(169|272|272|850|903|850\)&\(169|272|272\)\\
2&\(169|272|272\)&\(169|272|272|850|963|890\)&\(850|963|890\)\\
3&\(850|963|890\)&\(850|963|890|149|252|212\)&\(149|252|212\)\\
4&\(149|252|212\)&\(149|252|212|870|923|810\)&\(810|870|923\)\\
5&\(810|870|923\)&\(870|923|810|418|674|521\)&\(418|674|521\)\\
6&\(418|674|521\)&\(418|674|521|830|943|830\)&\(830|830|943\)\\
7&\(830|830|943\)&\(830|943|830|109|252|252\)&\(109|252|252\)\\
8&\(109|252|252\)&\(903|850|850|252|109|252\)&\(850|850|903\)\\
\hline
\end{tabular}
\end{center}

\paragraph{Testigo \(i=5\).}
\begin{center}
\small
\begin{tabular}{rlll}
\hline
\(j\)&\(v_{j-1}\)&\(e_j\)&\(v_j\)\\
\hline
1&\(850|850|903\)&\(109|252|252|850|903|850\)&\(109|252|252\)\\
2&\(109|252|252\)&\(109|252|252|501|614|498\)&\(498|501|614\)\\
3&\(498|501|614\)&\(501|614|498|169|272|272\)&\(169|272|272\)\\
4&\(169|272|272\)&\(169|272|272|850|963|890\)&\(850|963|890\)\\
5&\(850|963|890\)&\(850|963|890|149|252|212\)&\(149|252|212\)\\
6&\(149|252|212\)&\(149|252|212|830|943|830\)&\(830|830|943\)\\
7&\(830|830|943\)&\(830|943|830|109|252|252\)&\(109|252|252\)\\
8&\(109|252|252\)&\(903|850|850|252|109|252\)&\(850|850|903\)\\
\hline
\end{tabular}
\end{center}

\begin{theorem}[Minimalidad de las cinco coberturas regionales]
\label{regcl27:minimalidad}
Para cada \(i\in\{1,\ldots,5\}\), entre los paseos cerrados de \(G\)
que recorren al menos una vez cada arista de
\[
 E_i^*=\{U_{\rm APP},U_e,U_\varphi,U_\pi^{(i)}\},
\]
la longitud mínima, contada con multiplicidad de recorridos, es ocho.
Se permiten repeticiones de vértices y aristas.
\end{theorem}

\begin{proof}
Todo cierre que recorre \(U_{\rm APP}\) pasa por \(v_{\rm APP}\).
La rotación cíclica de su lista de pasos permite comenzar en ese vértice
sin cambiar longitud ni cobertura. Basta, por tanto, resolver el problema
con origen fijo.

El espacio de búsqueda es el conjunto finito
\[
 \mathscr S_i=V(G)\times\mathcal P(E_i^*),
 \qquad |\mathscr S_i|\le32\cdot16=512.
\]
El estado inicial es \(s_0=(v_{\rm APP},\varnothing)\), y el estado
terminal es \(s_*=(v_{\rm APP},E_i^*)\). Al recorrer una arista \(e\)
incidente en \(v\), hacia su otro extremo \(v'\), se actualiza
\begin{equation}
 (v,M)\longmapsto\bigl(v',M\cup(E_i^*\cap\{e\})\bigr).
 \label{regcl27:transicion-bfs}
\end{equation}
La máscara registra etiquetas de arista, no sólo extremos; de ahí que
no deban colapsarse aristas paralelas distintas.

Todas las transiciones cuestan una unidad. Una cola FIFO, iniciada en
\(s_0\), extrae estados por distancia creciente, marca cada estado en
su primera visita y guarda su predecesor y la etiqueta recorrida.
Equivalentemente, si \(\mathcal F_0=\{s_0\}\), las capas nuevas son
\[
 \mathcal F_{k+1}
 =\operatorname{Succ}(\mathcal F_k)
   \setminus\bigcup_{j=0}^{k}\mathcal F_j.
\]
Por inducción, \(\mathcal F_k\) contiene exactamente los estados cuya
distancia mínima desde \(s_0\) es \(k\). La enumeración finita sobre el
emisor indicado da, para cada una de las cinco regiones,
\[
 s_*\notin\bigcup_{k=0}^{7}\mathcal F_k,\qquad s_*\in\mathcal F_8,
 \qquad (\ell_1,\ell_2,\ell_3,\ell_4,\ell_5)=(8,8,8,8,8).
\]
Las tablas exhiben las cinco cotas superiores, incluyendo las cuatro
etiquetas requeridas en cada caso. La exploración exhaustiva de las
capas anteriores prueba la cota inferior. La combinación prueba la
minimalidad en el problema declarado.
\end{proof}

La longitud ocho caracteriza la cobertura conjunta de cuatro incidencias
en cada región de clausura. El cociente retiene las etiquetas de emisión
y sus empalmes entre clases de fase; la evolución temporal enriquecida
conserva, además, orientación y memoria. Esta descripción enlaza las cinco
regiones con los mismos anclajes de propagación, autoescala y APP, y
complementa el ciclo local de seis empalmes. La prolongación coinductiva
mantiene sus operadores y cilindros junto con esos datos de orientación
y memoria.
% END R28_REGIONAL_CLOSURES
